\documentclass[10pt,a4paper]{amsart}
\usepackage[margin=19mm]{geometry}
\usepackage{amsmath,amssymb,mathtools}
\usepackage[scr=rsfs,cal=euler]{mathalfa}
\usepackage{xfrac}
\usepackage[unicode,hidelinks,bookmarks=false]{hyperref}
\newcommand{\ie}{i.e.\ }
\newcommand{\cf}{cf.\ }
\newcommand{\resp}{respectively\ }
\newcommand{\Subsection}{\subsection}
\providecommand{\nobreakdash}{\nobreak\hbox{-}}
\setlength{\emergencystretch}{2em}
\multlinegap0pt
\usepackage{bm}
\usepackage{bbm}
\usepackage{dsfont}
\usepackage{xspace}
\usepackage{cancel}
\usepackage{mathtools}
\usepackage{cleveref}
\usepackage{amsthm}
\makeatletter
\@ifundefined{lemma}{\newtheorem{lemma}{Lemma}}{}
\makeatother
\DeclareMathOperator{\rank}{rk}
\title[Polar Degrees of Matroids]{Polar Degrees of Matroids}
\author{Clara Briand, Leonie Kayser, Julian Weigert}
\date{Source: arXiv:2606.26281v1 (CC BY 4.0)}
\begin{document}
\maketitle

\noindent\emph{Source and licence.} Polar Degrees of Matroids. Version arXiv:2606.26281v1 (CC BY 4.0), distributed under the Creative Commons Attribution 4.0 International licence, as recorded in the arXiv licence metadata for this exact version and shown on the abstract page https://arxiv.org/abs/2606.26281v1. Full rights notice: licenses/arxiv-2606.26281-CC-BY-4.0.txt.\par\smallskip

\noindent\textit{Editorial scope.} Counted unit: the Lemma whose source label is \texttt{lem: polynomialFormula} in the article (source0.tex lines 1038--1043, source label \texttt{lem: polynomialFormula}), delivered together with the article's own complete original proof of it (source0.tex lines 1044--1072, one \texttt{proof} environment of 29 lines). No mathematics is written, repaired or re-derived: statement and proof are the article's own text line for line. Required context retained uncounted: lem: formulaforqM (lines 1074--1076). Every definition, display and label of those lines is retained unchanged. Omitted from this package and not redistributed: 1--1037, 1073--1073, 1077--1504. Nothing inside a retained excerpt is omitted or altered: every retained body is delivered exactly as the article's lines are written, so no omission receipt of any kind accompanies this package. Citations: the retained excerpts carry no citation command at all, so no bibliography entry of the article is transcribed and no reference is redistributed. Reference adaptation: none was needed and none was made; every reference inside the delivered unit resolves inside it, so no unresolved reference or citation is produced. Printed slips kept literal: no printed word, symbol, hypothesis or number of the counted statement or of its proof was altered. Layout and class: the article's own class is \texttt{scrartcl} with packages including lmodern, babel, csquotes, xpatch, biblatex, enumitem, amsmath, mathtools, amssymb, amsthm, thmtools, tikz-cd, xcolor, braket; the wrapper is the lane's pinned \texttt{amsart} editorial wrapper at 10pt on A4, which restates the theorem-like environments the retained text needs and adds the article's own macro definitions that the retained text needs (\texttt{\textbackslash rank}), with extra wrapper packages (bm, bbm, dsfont, xspace, cancel, mathtools, cleveref). The delivered numbering is the unit's own. Assets: no retained span contains a figure environment, a graphics-inclusion command, a PDF-page inclusion, a table or a TikZ picture; no image file is copied into this package, the package declares no asset dependency and no image file is redistributed with it. Licence: version arXiv:2606.26281v1 of this article is distributed under the Creative Commons Attribution 4.0 International licence, as recorded in the arXiv licence metadata for this exact version and shown on the abstract page https://arxiv.org/abs/2606.26281v1. A full rights notice is delivered at licenses/arxiv-2606.26281-CC-BY-4.0.txt. Characters: every retained excerpt is pure ASCII, so the delivered engine, XeLaTeX with its default Latin Modern text font, reports no missing character.
\begin{lemma}\label{lem: formulaforqM}
In the notation above we have $q_M(t) = \overline{\chi}_M(t)$.
\end{lemma}

\begin{lemma}\label{lem: polynomialFormula}
We have the equality
\[
P_M(t) = (-t-1)^{d}\overline{\chi}\left(\frac{1}{t+1}\right)
\]
\end{lemma}

\begin{proof}
We proceed by induction on the rank of $M$. The base case is
$\rank(M) = 1$. In this case,
the only decreasing flag has length 0, so $P_M(t) = 1$, which agrees with the right hand side of the statement since the reduced characteristic polynomial is
$
\overline{\chi}_M(t) = 1
$.

Now let $\rank(M)>1$.
For each summand in $P_M(t)$ corresponding to a flag $\mathcal{F}$ of length $k>0$ we single out the largest proper non-empty flat $F=F_k$ of $\mathcal{F}$. We rewrite the summation defining $P_M(t)$ as a sum over flats $F$ appearing in this way. Notice that by the condition of being decreasing, we always have $0\notin F$. Including the summand $\beta(M)$ for the unique length 0 flag gives us the recursion 
\begin{align*}
    P_M(t)&=\beta(M)+\sum_{\substack{F \text{ flat} \\ 0\notin F\\F\neq \emptyset}}\sum_{\substack{\mathcal{F}' \text{ decr.} \\ \text{ flag in }M|_F}}\beta(M/F)\beta(M|_F[\mathcal{F}'])t^{|\mathcal{F}'|+1} \\
    &=\beta(M)+t\sum_{\substack{F \text{ flat} \\ 0\notin F\\F\neq \emptyset}}\beta(M/F)P_{M|_F}(t)
\end{align*}
Notice that every occurrence $P_{M|_F}(t)$ on the right hand side uses a strictly smaller matroid since we know $0\notin F$. Hence by induction we know \begin{align*}
        P_M(t)=\beta(M)+t\sum_{\substack{F \text{ flat} \\ 0\notin F\\F\neq \emptyset}}\beta(M/F)(-t-1)^{\rank_M(F)-1}\overline{\chi}_{M|_F}\big(\frac{1}{1+t}\big)
\end{align*}
We substitute $t$ by $\frac{1}{t}-1$ and multiply by $(-t)^{\rank(M)-1}$ to get \begin{align*}
    &(-t)^{\rank(M)-1}P_M\big(\frac{1}{t}-1\big)\\&=\beta(M)(-t)^{\rank(M)-1}+(t-1)\sum_{\substack{F \text{ flat} \\ 0\notin F\\F\neq \emptyset}}\beta(M/F)(-t)^{\rank(M)-\rank_M(F)-1}\overline{\chi}_{M|_F}(t) \\
    &=\beta(M)(-t)^{\rank(M)-1}+\sum_{\substack{F \text{ flat} \\ 0\notin F\\F\neq \emptyset}}\beta(M/F)(-t)^{\rank(M)-\rank_M(F)-1}\chi_{M|_F}(t) \\
        &=\sum_{\substack{F \text{ flat} \\ 0\notin F}}\beta(M/F)(-t)^{\rank(M)-\rank_M(F)-1}\chi_{M|_F}(t) \\
\end{align*}
Here in the last step we include again the summand for the empty flat using that $\chi_{M|_\emptyset}(t)=1$. We could not include this summand earlier since it is the only summand where the non-reduced characteristic polynomial is not divisible by $t-1$, hence this summand has no reduced characteristic polynomial. 

By applying the same substitutions to the equality that we want to prove, we see that it remains to show \begin{align*}
    \overline{\chi}_M(t)=\sum_{\substack{F \text{ flat} \\ 0\notin F}}\beta(M/F)(-t)^{\rank(M)-\rank_M(F)-1}\chi_{M|_F}(t) \eqcolon q_M(t)
\end{align*}
which is done in the following \Cref{lem: formulaforqM}.
\end{proof}

\end{document}
